\section{Cronograma de actividades}

El trabajo se acoge a la \textbf{opción 2} del calendario, sin opción a mención de honor y con ceremonia de grados el 9 y 10 de diciembre de 2026. El período 2026-2 inicia en la semana del 13 de julio (semana~1), incluye una semana de receso y cierra con el balance académico el 25 de noviembre. El hito que gobierna la planeación no es la defensa sino la entrega del informe final con carta de aval el 30 de octubre (semana~15): toda la experimentación debe cerrarse en la semana~14. La redacción del informe no se trata como una etapa final, sino como una actividad transversal que arranca en la semana~3 —con el anteproyecto ya escrito— y se alimenta de cada fase a medida que se ejecuta, de modo que las últimas quincenas se dediquen a consolidar y revisar, no a escribir desde cero. El cronograma de actividades se muestra en la figura~\ref{fig:cronograma}.  

\begin{figure}[H]
\centering
\includegraphics[width=\linewidth]{cronograma}
\caption{Cronograma del período académico 2026-2.}
\label{fig:cronograma}
\end{figure}
